\section{The lowest coefficient at a scalar triple}
\label{sch:sec:scalar}

At a scalar triple, every subspace is invariant. Its Hall coefficient
still depends on the ambient cubic potential. Complete flags produce
a distinguished class in this coefficient, including at the scalar point.
The class determines the lowest-degree images for the potentials of
ranks three, four, and five.

All spaces are complex algebraic varieties with their analytic topology.
Coefficients are rational. For $1\leq n\leq5$, put
\[
 V_n=\operatorname{Mat}_n(\mathbb C)^3,\qquad
 f_n(A,B,C)=\operatorname{Tr}(A[B,C]),\qquad
 Z_n=\{[A,B]=[B,C]=[C,A]=0\}.
\]
Differentiation using the trace pairing gives $Z_n=\operatorname{Crit}(f_n)$.
Write $\mathfrak X_n=[V_n/\operatorname{GL}_n]$ and
$\mathfrak M_n=[Z_n/\operatorname{GL}_n]$. Constructible complexes on
these quotient stacks retain the conjugation action and its descent data.
Their ordinary atlas pullbacks have no additional shift. The coefficient is
\[
 \mathcal P_n=
 \left.\Phi_{f_n}\mathbb Q_{\mathfrak X_n}[2n^2]\right|_{\mathfrak M_n},
 \qquad
 \Phi_fK=\operatorname{Cone}(i^*K\longrightarrow\psi_fK)[-1],
\]
where $i:f^{-1}(0)\hookrightarrow V$ and $\psi_f$ denotes nearby cycles.
This normalization agrees with Massey\cite[pp.~1--2]{schMasseyST}.
Thus $H^j(K[a])=H^{j+a}(K)$ and $\Phi_0K=K$.
The tensor differential is
$d(x\otimes y)=dx\otimes y+(-1)^{|x|}x\otimes dy$.
No completion or change of coefficients is used.

For $s=(aI_n,bI_n,cI_n)$, write $P_{n,s}$ for the ordinary atlas
stalk of $\mathcal P_n$. Translation by scalar matrices preserves $f_n$,
so it suffices to compute $P_{n,0}$.
Define the commuting-pair variety and its complement by
\[
 \mathcal C_n=\{(A,C)\in\operatorname{Mat}_n(\mathbb C)^2:[C,A]=0\},
 \qquad U_n=\operatorname{Mat}_n(\mathbb C)^2\setminus\mathcal C_n.
\]
Borel--Moore homology below uses locally finite rational chains and
the complex orientation of each smooth stratum.

\begin{proposition}[Scalar stalk and commuting pairs]
\label{sch:prop:scalar}
For $2\leq n\leq5$ and every integer $j$, there are isomorphisms
\begin{equation}
 H^j(P_{n,s})\simeq
 \widetilde H^{j+2n^2-1}(U_n;\mathbb Q)
 \simeq H^{\mathrm{BM}}_{2n^2-j}(\mathcal C_n;\mathbb Q).
 \label{sch:eq:scalar-bm}
\end{equation}
In particular,
\[
 H^j(P_{n,s})=0\quad(j<-2n),\qquad
 H^{-2n}(P_{n,s})\simeq\mathbb Q.
\]
The same last two assertions hold for $n=1$, where
$P_{1,s}=\mathbb Q[2]$. Potential monodromy acts as the identity on
the entire scalar stalk complex for every rank under consideration.
\end{proposition}

\begin{proof}
Use the positive real ray in the potential plane. By
Massey\cite[\S1, p.~2]{schMasseyST}, the normalized
vanishing-cycle stalk is computed by cochains with support in
$P=\{\operatorname{Re}f_n\leq0\}$, restricted to the origin.
The fibre convention is the relative-cochain convention
$\operatorname{Fib}(C^*(B)\to C^*(B\setminus P))$.
Both $P$ and its complement are invariant under positive real scaling.
The radial map
\[
 x\longmapsto \frac{x}{1+\|x\|/\delta}
\]
maps $V_n$ into the ball of radius $\delta$. Interpolation of the
denominator gives homotopy inverses for the inclusions of these pairs.
Thus their relative cohomology computes the scalar stalk as
$\delta$ decreases to zero.

Let $F_n=f_n^{-1}(1)$ and $H=\{z\in\mathbb C:\operatorname{Re}z>0\}$.
Choose the cubic root on $H$ that is positive on the positive real axis.
Homogeneity gives mutually inverse homeomorphisms
\[
 V_n\setminus P\longrightarrow F_n\times H,
 \quad x\longmapsto(f_n(x)^{-1/3}x,f_n(x)),
 \qquad (y,z)\longmapsto z^{1/3}y.
\]
Since $V_n$ and $H$ are contractible, the support complex has
cohomology $\widetilde H^{k-1}(F_n;\mathbb Q)$ in degree $k$.
This argument fixes the local-to-global comparison at the scalar point.

The cyclic trace identity writes the actual potential as
$f_n=\operatorname{Tr}(B[C,A])$. Projection onto $(A,C)$ gives
\begin{equation}
 F_n\longrightarrow U_n,
 \qquad (A,B,C)\longmapsto(A,C),
 \label{sch:eq:affine-fibre}
\end{equation}
an affine bundle of complex rank $n^2-1$. Indeed, its fibre is the
affine hyperplane $\operatorname{Tr}(BD)=1$, where $D=[C,A]\ne0$.
The trace pairing is nondegenerate. A continuous section is
\[
 B_0(A,C)=\frac{D^*}{\operatorname{Tr}(D^*D)},
\]
where $D^*$ is the conjugate transpose. The affine homotopy from $B$
to $B_0$ stays in this hyperplane. Therefore projection induces the
first isomorphism in \eqref{sch:eq:scalar-bm}, after the shift $[2n^2]$.
Projection is conjugation equivariant. The chosen section only supplies
an ordinary topological homotopy inverse.

The real dimension of the pair space is $4n^2$. Its only nonzero
Borel--Moore group is $H^{\mathrm{BM}}_{4n^2}=\mathbb Q$.
The closed-open exact sequence for $\mathcal C_n$ and $U_n$, followed
by Poincar\'e duality on the smooth open set $U_n$, gives
\[
 \widetilde H^k(U_n;\mathbb Q)
       \simeq H^{\mathrm{BM}}_{4n^2-k-1}(\mathcal C_n;\mathbb Q).
\]
For $k=0$, the reduced group removes the ambient fundamental class.
The complement of a proper complex subvariety in this affine space
is connected: a generic complex line joins each point to a common
generic intermediate point while avoiding finitely many intersections.
The remaining nonpositive degrees vanish. Substitution gives the
second isomorphism in \eqref{sch:eq:scalar-bm}.

For completeness, the top dimension of $\mathcal C_n$ follows directly
from Jordan forms. Fix a Jordan type for $A$ with $r$ distinct eigenvalues
and a centralizer of dimension $c$. Its matrix stratum has dimension
$n^2-c+r$. The matrices $C$ commuting with $A$ form a vector bundle
of rank $c$ on this stratum. The corresponding pair stratum has
dimension $n^2+r$.
There are finitely many Jordan types. The unique stratum with $r=n$
has dimension $n^2+n$ and is irreducible: ordered eigenvalues and
diagonal entries of $C$ parameterize it, together with a change of basis.
The finite permutation quotient preserves irreducibility.
Every other stratum has dimension at most $n^2+n-1$.
Consequently the top Borel--Moore group is
\[
 H^{\mathrm{BM}}_{2n^2+2n}(\mathcal C_n;\mathbb Q)=\mathbb Q,
 \qquad H^{\mathrm{BM}}_m(\mathcal C_n;\mathbb Q)=0
                 \quad(m>2n^2+2n).
\]
The top generator is the complex fundamental class of the closure
of that stratum. No assertion about lower-dimensional components
is needed. Equation~\eqref{sch:eq:scalar-bm} gives the degree bounds.

Finally, at the scalar origin the maps
\[
 (A,B,C)\longmapsto(A,e^{i\theta}B,C),\qquad0\leq\theta\leq2\pi,
\]
preserve every centered Euclidean ball and multiply $f_n$ by
$e^{i\theta}$. They trivialize the Milnor fibres over each circle
in the potential plane. At $2\pi$ the map is the identity.
These trivializations commute with restriction to smaller balls,
with the constant-to-nearby-cycle map, and with conjugation.
They give identity potential monodromy on the scalar fibre complex.
Translation gives the same assertion at every scalar triple.
\end{proof}

The first nonzero groups for the three target potentials are therefore
\[
\begin{array}{c|c|c|c}
 n&\dim_{\mathbb C}\mathcal C_n&
  \text{first reduced Milnor degree}&\text{first degree of }P_{n,s}\\\hline
 3&12&11&-6\\
 4&20&23&-8\\
 5&30&39&-10.
\end{array}
\]
Each displayed group has rank one. The calculation concerns ordinary
atlas stalks. It does not replace equivariant coefficient complexes
by these graded vector spaces.

The rank-two scalar fibre gives a further check on the shifts.
Write the traceless parts of $A,B,C$ in the basis
$\operatorname{diag}(1,-1),e_{12},e_{21}$ of
$\mathfrak{sl}_2$. If $T$ is the matrix of their three coordinate
columns, direct expansion gives $f_2=2\det T$.
The scalar directions are free. The determinant-one fibre retracts
onto $\operatorname{SU}(3)$ by polar decomposition.
The fibration $\operatorname{SU}(2)\to\operatorname{SU}(3)\to S^5$
has fibre $S^3$. Its rational Serre spectral sequence has nonzero
terms only in bidegrees $(0,0),(0,3),(5,0),(5,3)$, so every
differential vanishes. The reduced fibre cohomology has degrees
$3,5,8$, each of rank one. After the normalized shift $[8]$,
the scalar coefficient has degrees $-4,-2,1$.

\section{Complete flags define scalar generators}
\label{sch:sec:flags}

The rank-one coefficient is $\mathbb Q[2]$. Complete flags let these
coefficients enter the higher-rank critical coefficient without a choice
of a scalar Milnor cycle.
Let $\operatorname{Fl}_n$ be the variety of complete flags in
$\mathbb C^n$ and set $p_n=n(n-1)/2$.
Let $E_n^{\mathrm{fl}}\subset V_n\times\operatorname{Fl}_n$ classify
triples preserving a complete flag. Its matrices are upper triangular
in an adapted basis, so $f_n|_{E_n^{\mathrm{fl}}}=0$.
The normal bundle of this incidence has complex rank $3p_n$.
Its positive complex Thom class followed by flag integration defines
\[
 \gamma_n^{\mathrm{fl}}:
 Rq_{n*}\mathbb Q_{E_n^{\mathrm{fl}}}
       \longrightarrow\mathbb Q_{V_n}[4p_n],
\]
where $q_n$ forgets the flag. The map is proper because
$\operatorname{Fl}_n$ is projective.

Pass to quotient stacks, and retain the full critical fibre product
\[
 \mathcal T_n=[E_n^{\mathrm{fl}}/\operatorname{GL}_n]
                    \mathbin{\times}_{\mathfrak X_n}\mathfrak M_n,
 \qquad b_n:\mathcal T_n\longrightarrow\mathfrak M_n.
\]
Thus the off-diagonal commutator equations remain part of the source.
Proper exchange for vanishing cycles
\cite[\S3, p.~35]{schMasseyNotes} and $\Phi_0=1$ give
\[
 \nu_n:Rb_{n*}\mathbb Q_{\mathcal T_n}[2n]
       \longrightarrow\mathcal P_n,
 \qquad 2n+4p_n=2n^2.
\]
The unit of pullback and pushforward then defines degree-zero morphisms
\begin{equation}
 h_n:\mathbb Q_{\mathfrak M_n}[2n]
 \longrightarrow Rb_{n*}\mathbb Q_{\mathcal T_n}[2n]
 \xrightarrow{\nu_n}\mathcal P_n,
 \qquad \sigma_n=\frac{h_n}{n!}.
 \label{sch:eq:class}
\end{equation}
Every construction commutes with changes of frame. Its normal bundles
and flag cycles retain their complex orientations on the action diagram.
Thus \eqref{sch:eq:class} defines morphisms on the quotient stacks.

The sign of $h_n$ can be determined at a matrix $A$ with distinct
eigenvalues $\lambda_1,\ldots,\lambda_n$. In a transverse eigenbasis
chart, the actual summand for $i<j$ is
\[
 (\lambda_i-\lambda_j)(b_{ij}c_{ji}-b_{ji}c_{ij}).
\]
Set $\delta_{ij}=\lambda_i-\lambda_j\ne0$ and introduce
\[
 \widehat c_{ij}=\delta_{ij}c_{ij},\qquad
 \widehat c_{ji}=\delta_{ij}c_{ji},\qquad
 q=b_{ij}\widehat c_{ji}-b_{ji}\widehat c_{ij}.
\]
Each rescaling is complex linear with real determinant
$|\delta_{ij}|^2>0$. Both together have determinant
$|\delta_{ij}|^4>0$, so they preserve the complex normal orientation.
The $B$-polarized positive disk has normalized coordinates
\[
 (b_{ij},b_{ji},\widehat c_{ij},\widehat c_{ji})=(u,v,-\bar v,\bar u),
 \qquad q=|u|^2+|v|^2.
\]
Its complex parameter orientation defines a positive support generator.
Projection to the normalized lower normal coordinates is
$(u,v)\mapsto(v,\bar u)$, of real degree $-1$. In the original
coordinates it is $(u,v)\mapsto(v,\delta_{ij}^{-1}\bar u)$.
The additional complex multiplication has positive real determinant
$|\delta_{ij}|^{-2}$, so the projection degree remains $-1$.
The lower $A$ normal cancels the flag tangent by multiplication with
the nonzero eigenvalue difference. This complex linear map has
positive real determinant.
If $u_n$ denotes the product of the $B$-polarized generators, put
\[
 e_n=(-1)^{p_n}u_n.
\]
All support Thom degrees are even. Their ordered external products
introduce no sign. Thus each complete flag contributes $e_n$, and
there are $n!$ complete flags. This proves
\begin{equation}
 h_n|_{\mathrm{distinct}}=n!e_n,
 \qquad \sigma_n|_{\mathrm{distinct}}=e_n.
 \label{sch:eq:regular}
\end{equation}
Here $\mathcal P_n=\mathbb Q[2n]$ on this locus, since its normal
quadratic has $2(n^2-n)$ variables.

\begin{theorem}[Extension of the normalized lowest coefficient]
\label{sch:thm:generator}
For $1\leq n\leq5$, the stalk of $\sigma_n$ at any scalar triple
induces an isomorphism
\[
 \mathbb Q\xrightarrow{\ \sim\ }H^{-2n}(P_{n,s}).
\]
Denote the image of $1$ by $v_{n,s}$. On the scalar locus, these
classes form a constant, conjugation-equivariant cohomology sheaf.
They have identity potential monodromy.
\end{theorem}

\begin{proof}
Positive scaling of all three matrices preserves $Z_n$ and each
flag incidence, and multiplies $f_n$ by a positive cube.
It preserves the support half-plane, the flag cycles and the Thom
orientations. Hence the section of $\mathcal H^{-2n}(\mathcal P_n)$
induced by $h_n$ is equivariant under positive scaling.

Suppose this section had zero stalk at the origin on the atlas.
By the definition of a sheaf stalk, it would vanish on a neighborhood
of the origin. Every point of $Z_n$ can be scaled into that neighborhood.
Scaling equivariance would force the section to vanish everywhere.
This contradicts \eqref{sch:eq:regular}. Therefore $h_n$ has a
nonzero scalar stalk in degree $-2n$.
Proposition~\ref{sch:prop:scalar} shows that this stalk group is
one-dimensional. Division by $n!$ gives the stated isomorphism.

Scalar translations preserve the function, incidences and orientations.
They transport this generator to every scalar triple.
The morphism already descends under conjugation, so its lowest
cohomology sheaf has the trivial scalar stabilizer action.
The monodromy assertion also follows from
Proposition~\ref{sch:prop:scalar}.
\end{proof}

This theorem specifies the lowest cohomology sheaf. It supplies no
splitting of the higher cohomology or of the equivariant derived object.

\section{The scalar images of the Hall coefficient maps}
\label{sch:sec:images}

Fix positive integers $d,e$ with $n=d+e\leq5$. The product order is
submodule first, followed by quotient. Let $E_{d,e}$ classify triples
preserving a $d$-plane, and write
\[
 \mathcal F_{d,e}=[E_{d,e}/\operatorname{GL}_n]
                         \mathbin{\times}_{\mathfrak X_n}\mathfrak M_n,
 \qquad
 \mathfrak M_d\times\mathfrak M_e
 \xleftarrow{\ a_{d,e}\ }\mathcal F_{d,e}
 \xrightarrow{\ b_{d,e}\ }\mathfrak M_n.
\]
The diagonal block potential is $f_d+f_e$.
The normal Thom map has degree $6de$, and Grassmannian integration
has degree $-2de$. Their composite is
\[
 \gamma_{d,e}:Rq_{d,e*}\mathbb Q_{E_{d,e}}
                      \longrightarrow\mathbb Q_{V_n}[4de].
\]
For analytic functions $f,g$, the relative-cochain product and the
support inclusion
\[
 \{\operatorname{Re}f\leq0\}\times\{\operatorname{Re}g\leq0\}
       \subset\{\operatorname{Re}(f+g)\leq0\}
\]
define the morphism
$\Phi_fK\boxtimes\Phi_gL\to k^*\Phi_{f+g}(K\boxtimes L)$.
The product and support inclusion are those of
Massey\cite[Lemma~1.2, p.~3, and Proposition~1.4, p.~5]{schMasseyST}.
Here $k$ includes the product of zero fibres in the zero fibre of the sum.
This morphism is natural in coefficients and proper direct image:
both assertions follow from the same maps of relative cochain complexes.
For proper direct image, apply proper base change and the proper
K\"unneth isomorphism\cite[\S1, pp.~13--14]{schMasseyNotes}.
Its ordered triple products agree because they use the same product
cochains and the same inclusion into the support of $f+g+h$.

Use this support morphism, smooth ambient block pullback, proper
exchange, and $\Phi(\gamma_{d,e})$. Restriction to the full critical
fibre product defines the actual coefficient map
\begin{equation}
 \mu_{d,e}:Rb_{d,e*}a_{d,e}^*
       (\mathcal P_d\boxtimes\mathcal P_e)\longrightarrow\mathcal P_n.
 \label{sch:eq:mu}
\end{equation}
Its shift identity is $2d^2+2e^2+4de=2n^2$.
The critical block map need not be smooth.
The construction uses the support comparison as a morphism, including
for $f_2+f_3$. Its invertibility is unnecessary here.

Let $\eta_{d,e}$ denote the unit
$\mathbb Q_{\mathfrak M_n}\to Rb_{d,e*}\mathbb Q_{\mathcal F_{d,e}}$.
The canonical tensor identification and the maps $\sigma_d,\sigma_e$
give a morphism from $\mathbb Q_{\mathcal F_{d,e}}[2n]$ to the
incoming coefficient in \eqref{sch:eq:mu}.

\begin{theorem}[Scalar lowest-degree Hall images]
\label{sch:thm:images}
For $d+e=n\leq5$, the following equality holds as morphisms of
constructible complexes on $\mathfrak M_n$:
\begin{equation}
 \mu_{d,e}\circ Rb_{d,e*}a_{d,e}^*
       (\sigma_d\boxtimes\sigma_e)\circ\eta_{d,e}[2n]
       =\binom nd\sigma_n.
 \label{sch:eq:divided}
\end{equation}
At a scalar triple, the source and target of the induced map in
degree $-2n$ are both one-dimensional. In the generators specified
by $v_{d,s}$, $v_{e,s}$ and $v_{n,s}$, this map is multiplication
by $\binom nd$. Its image is the whole lowest-degree target and
its kernel in that degree is zero.
\end{theorem}

\begin{proof}
Refine an invariant $d$-plane by a complete flag in that plane and
a complete flag in its quotient. The result is exactly a complete
flag in $\mathbb C^n$. This identifies the refinement square on
the full critical stacks, including its frame changes.
Proper base change identifies the corresponding iterated direct images.

The forbidden lower entries for the refined flag are the disjoint
union of internal entries in the two blocks and the cross entries.
The product of their complex Thom classes is the complete-flag Thom
class. Integration over the two internal flag varieties and then
the Grassmannian evaluates on the same complete-flag fundamental cycle.
The identity $p_n=p_d+p_e+de$ gives the total degree $4p_n$.
All these degrees are even.

On coefficients, proper naturality of the support product commutes
the internal $\Phi(\gamma_d^{\mathrm{fl}})$ and
$\Phi(\gamma_e^{\mathrm{fl}})$ through the external support morphism.
Before the internal pushforwards, all one-dimensional block potentials
are zero. Their ordered support product is the constant coefficient
on the complete-flag incidence. After the pushforwards, the target
potentials remain the actual $f_d$ and $f_e$.
The composite therefore equals $h_n$ when its inputs are $h_d,h_e$.
The adjunction units compose to the complete-flag unit because they
send the constant section to the same constant section on the refinement.
Dividing by $d!e!$ and using $h_n=n!\sigma_n$ proves
\eqref{sch:eq:divided}.

At a scalar atlas point, proper base change identifies the incoming
stalk with hypercohomology on $\operatorname{Gr}(d,n)$ of the
pulled-back coefficient family. Proposition~\ref{sch:prop:scalar}
puts that family in degrees at least $-2d-2e=-2n$.
Theorem~\ref{sch:thm:generator} identifies its cohomology sheaf in
degree $-2n$ with the constant sheaf generated by $v_d\otimes v_e$.
The hypercohomology spectral sequence has no term below this degree.
Its degree $-2n$ term is
$H^0(\operatorname{Gr}(d,n),\mathbb Q)=\mathbb Q$.
It has no incoming or outgoing differential, since the other relevant
terms would have negative sheaf-cohomology degree or lower coefficient
degree. Equation~\eqref{sch:eq:divided} proves the scalar assertion.
\end{proof}

In particular the first images involving $f_3,f_4,f_5$ are
\[
\begin{array}{c|c|c|c}
 (d,e)&\text{degree}&\text{incoming lowest class}&\text{image}\\\hline
 (2,1)&-6&v_2\otimes v_1&3v_3\\
 (2,2)&-8&v_2\otimes v_2&6v_4\\
 (3,1)&-8&v_3\otimes v_1&4v_4\\
 (4,1)&-10&v_4\otimes v_1&5v_5\\
 (2,3)&-10&v_2\otimes v_3&10v_5.
\end{array}
\]
Each incoming class denotes the constant section on the corresponding
Grassmannian, with the coefficient morphisms just constructed.

\section{Three constituents and the remaining scalar coefficients}
\label{sch:sec:coherence}

For an ordered triple $(d,e,f)$ of positive ranks with sum $n\leq5$,
let $\mathcal F_{d,e,f}$ be the full critical flag stack with these
successive dimensions. Its maps $a$ and $b$ remember the constituents
and the total object. The ordered triple support product and the flag
Thom map define
\[
 M_{d,e,f}:Rb_*a^*(\mathcal P_d\boxtimes\mathcal P_e
                         \boxtimes\mathcal P_f)\longrightarrow\mathcal P_n.
\]
Both parenthesizations use this source, by proper base change.
Their ambient Gysin degrees add to
$4(de+df+ef)$, since
\[
 de+(d+e)f=ef+d(e+f).
\]
Their normal entries partition the same forbidden entries, and their
flag integrations evaluate on the same complex cycle. Ordered support
associativity and its coefficient naturality therefore give equality
of the two derived maps with $M_{d,e,f}$.
This argument retains $f_{d+e}$ and $f_{e+f}$ on the intermediate
representation spaces. It introduces no permutation of the inputs.

Refinement to complete flags gives the following restriction of this
identity to the constructed coefficient morphisms:
\begin{equation}
 M_{d,e,f}(\sigma_d,\sigma_e,\sigma_f)
                  =\frac{n!}{d!e!f!}\sigma_n.
 \label{sch:eq:triple}
\end{equation}
The left side includes the adjunction unit, as in
\eqref{sch:eq:divided}. At a scalar point its source in degree $-2n$
is again one-dimensional, by the same hypercohomology argument on
the connected flag variety.
For the first three scalar comparisons the two factorizations give
\[
\begin{array}{c|c|c|c}
 (d,e,f)&\text{left factorization}&\text{right factorization}&\text{image}\\\hline
 (1,1,1)&2\cdot3&2\cdot3&6v_3\\
 (2,1,1)&3\cdot4&2\cdot6&12v_4\\
 (2,2,1)&6\cdot5&3\cdot10&30v_5.
\end{array}
\]
These are images of scalar coefficient classes. They follow from
the complete-flag morphism and its nonzero scalar restriction.

The lowest coefficient has support on the entire scalar locus
$\mathbb C^3\times B\operatorname{GL}_n$ and on the distinct-eigenvalue
locus. The translation and conjugation compatibilities specify the
constant lowest cohomology sheaf on the scalar locus. The ambient
derived coefficient remains in every incidence and refinement.

Higher scalar images require further calculations. For example, the
rank-two scalar coefficient has degrees $-4,-2,1$. On a scalar
$(2,1)$ incidence, its degree $-2$ coefficient tensored with the
line coefficient starts in total degree $-4$.
That degree also receives $H^2(\mathbb P^2,\mathbb Q)$ from the
lowest coefficient. The morphism from these contributions to
$H^{-4}(P_{3,s})$ is not determined by
\eqref{sch:eq:divided}. It requires the actual coefficient family,
its hypercohomology differentials, and its map to the scalar target.
Likewise, the odd rank-two class remains part of the $(2,2)$ and
$(2,2,1)$ inputs. Equation~\eqref{sch:eq:scalar-bm} describes the
remaining scalar target groups through commuting-pair homology.
It does not identify the Hall maps on those groups.

The constructed scalar images and derived flag identities are rational
constructible statements through rank five. A comparison with a
holomorphic field-theory carrier requires an additional morphism and
compatibility with its operations. Tate structures, higher homotopies
and further ranks require their own constructions.
